\section{Renormalization of gluon-gauge-link vertex}
\label{sec:1g-link}
For the definiteness of following discussion, we assume that the gauge link in diagrams of the topology (c) in Fig.~\ref{fig:div-topo-g} starts at an arbitrary coordinate $\xi_{1z}$ with the operator $F^{\mu \nu}(\xi_{1z})$, and ends at another arbitrary coordinate $\xi_{2z}$ with no additional operators. With the ``bare'' coupling constant $g_s$, and ``bare'' field operators for the gluons, the Faddeev-Popov ghost and the antighost given by the symbols $A$, $c$ and $\bar{c}$ respectively, a generalized Ward identity of the non-Abelian field relevant to this topology can be derived \cite{Collins:1984xc},
\begin{align}\label{eq:gwi1}
\begin{split}
\langle -i \partial_\lambda^y A_d^\lambda(y)[\Phi(\{\xi_{2z},\xi_{1z}\})]_{a b} \,  F_{b}^{\mu \nu}(\xi_{1z}) \rangle\\
=\langle g_s \bar{c}_d (y) c_e (\xi_{2z}) [ t_e \Phi(\{\xi_{2z},\xi_{1z}\})]_{a b} \, F_{b}^{\mu \nu}(\xi_{1z}) \rangle ,
\end{split}
\end{align}
where the $t$ represents SU(3) generators of the adjoint representation.
A pictorial representation of Eq.~(\ref{eq:gwi1}) is given in Fig.~\ref{fig:gwi1}, where ``1PR'' denotes one-particle-reducible diagrams. The topology of the left-hand side of Fig.~\ref{fig:gwi1} is the same as that of the diagram (c) in Fig.~\ref{fig:div-topo-g}, but is contracted with external gluon momentum and expressed in coordinate space. The topology of the first term on the right-hand side of Fig.~\ref{fig:gwi1} is nonlocal in spacetime, and thus has no UV divergence. Furthermore, after the renormalization of QCD Lagrangian and gauge-link-related topologies (a) and (b) in Fig.~\ref{fig:div-topo-g}, UV divergences of 1PR diagrams will be canceled. That is, the generalized Ward identity in Eq.~(\ref{eq:gwi1}) ensures that the topology (c) in Fig.~\ref{fig:div-topo-g} is free of UV divergence if it is contracted by external gluon momentum.

%%%%%%%%%%%%%%%%%%%%%%%%%%%%%%%
\begin{figure}[htb]
	\begin{center}
		\includegraphics[width=0.55\textwidth]{images/Green-1-g.pdf}
		\caption{\label{fig:gwi1} Pictorial representation of the generalized Ward identity in Eq.~(\ref{eq:gwi1}) with dashed line represents the ghost field.}
	\end{center}
\end{figure}
%%%%%%%%%%%%%%%%%%%%%%%%%%%%%%

To understand the renormalization of the UV divergence of the topology (c) in Fig.~\ref{fig:div-topo-g}, we represent the diagrams of this topology as $\Gamma^{\lambda \mu \nu}(p,n)$, which could be referred as the gluon-gauge-link vertex, where $p$ and $\lambda$ are the momentum and Lorentz index of the external gluon, respectively. Lorentz symmetry combined with antisymmetry between $\mu$ and $\nu$ enable us to do the general decomposition $\Gamma^{\lambda \mu \nu}(p,n)=\sum_{i=1}^4 c_i \Pi_i^{\lambda \mu \nu}$ with
\begin{align}
\begin{split}
\Pi_1^{\lambda \mu \nu}&=g^{\mu\lambda} p^\nu-g^{\nu\lambda}p^\mu,\, \Pi_2^{\lambda \mu \nu}=(p^\mu n^\nu-p^\nu n^\mu)n^\lambda,\\
\Pi_3^{\lambda \mu \nu}&=(p^\mu n^\nu-p^\nu n^\mu)p^\lambda, \,\Pi_4^{\lambda \mu \nu}=g^{\mu\lambda} n^\nu-g^{\nu\lambda}n^\mu.
\end{split}
\end{align}
Since $p_\lambda \Gamma^{\lambda \mu \nu}=0$ for the UV divergent terms as discussed above, we obtain $c_2\, p\cdot n + c_3\, p^2 + c_4=0$, and consequently,
\begin{align}
\begin{split}
\Gamma^{\lambda \mu \nu}(p,n)=&c_1 \Pi_1^{\lambda \mu \nu} + c_2 (\Pi_2^{\lambda \mu \nu}-p\cdot n \Pi_4^{\lambda \mu \nu})\\
&+c_3 (\Pi_3^{\lambda \mu \nu}-p^2 \Pi_4^{\lambda \mu \nu}).
\end{split}
\end{align}
Since $c_1$ and $c_2$ have mass dimension $0$, locality of UV divergences ensures that they can be at most logarithmic divergent, while $c_3$ is UV finite due to its mass dimension at $-1$. The only potential linearly UV divergent coefficient $c_4$ is removed by gauge invariance. We have therefore demonstrated that the cancellation of linear UV divergences of diagrams (b) and (c) in Fig.~\ref{fig:oneloop-g} at one-loop order can be generalized to all orders, which makes the multiplicative renormalizability of quasi-gluon operators a possibility.

At the lowest order in $\alpha_s$, we have $\Gamma^{\lambda \mu \nu}(p,n) \propto \Pi_1^{\lambda \mu \nu}$. If we want $\Gamma^{\lambda \mu \nu}(p,n)$ not to mix with other operators under renormalization, we need its UV divergence to be proportional to $\Pi_1^{\lambda \mu \nu}$ to all orders. Fortunately, it is always true. For the case with $\mu$ (or $\nu$) along $z$-direction or the case with both $\mu$ and $\nu$ not along $z$-direction, the coefficients of $c_2$ are proportional to $\Pi_1^{\lambda \mu \nu}$ or equal to zero, respectively. Therefore, the components of $\Gamma^{\lambda \mu \nu}(p,n)$ do not mix with each others at all, although two different renormalization constants are needed for the two different choices.

In summary, if we choose either $F^{z \bar{\nu}}$ or $F^{\bar{\mu} \bar{\nu}}$ for gluon-gauge-link vertex, $\Gamma^{\lambda \mu \nu}$, we can remove the UV divergences of the vertex by multiplying a corresponding renormalization factor $Z_{vg1}^{-1/2}$ or $Z_{vg2}^{-1/2}$, respectively.
